\chapter{Scheduling and sequencing}
\label{chap:calendarizacion}

\section{Problem and assumptions}
Scheduling assigns operations to machines and times. QSB minimizes the
\emph{makespan}, the time at which the last operation finishes
\cite{pinedo2016scheduling}. Operations cannot be interrupted, a machine
processes one operation at a time, and each job follows its route. Durations
are integers in the specified unit. Setup times, working calendars, and
random failures are not included.
Due-date and weight fields are preserved, but the native objective remains
makespan; it does not optimize weighted tardiness.

\section{Flow shop: a common route}
In a flow shop, all jobs visit the machines in the same order.
The native solver searches for a permutation shared by all machines. For
a given sequence, an operation starts when both the job and machine are
available; its finish is its start plus its duration. Initial machine
availability times are respected when provided.

In \path{manual/examples/plan_flowShop.json}, there are two jobs and two
machines, initially available at zero:
\begin{center}
\begin{tabular}{lrr}\toprule
Job & M1 & M2 \\
\midrule J1 & 2 & 3 \\ J2 & 3 & 1 \\
\bottomrule\end{tabular}
\end{center}
\begin{lstlisting}[language=bash]
swift run qsb validate-scheduling-json manual/examples/plan_flowShop.json
swift run qsb solve-scheduling-json manual/examples/plan_flowShop.json --backend native
\end{lstlisting}
The sequence J1, J2 has makespan 6. On M1, J1 occupies $[0,2]$ and J2 $[2,5]$;
on M2, J1 occupies $[2,5]$ and J2 $[5,6]$. The alternative J2, J1 finishes at 8.
M2's initial idle time is not a missing operation: it is waiting for J1 to
finish. The search is exact within the permutation model, with a limit of
eight jobs and factorial cost; it does not prove optimality over all possible
machine-specific sequences in a general flow shop.

\section{Job shop: individual routes}
In a job shop, each job may follow a different route. The example
\path{plan_jobShop.json} requires J1: M1(2), M2(3), and J2: M2(2), M1(1).
Precedence constraints follow from the order of each job's operations.
\begin{lstlisting}[language=bash]
swift run qsb validate-scheduling-json manual/examples/plan_jobShop.json
swift run qsb solve-scheduling-json manual/examples/plan_jobShop.json --backend native
\end{lstlisting}
QSB obtains makespan 5: J1 uses M1 from 0 to 2 and M2 from 2 to 5; J2 uses
M2 from 0 to 2 and M1 from 2 to 3. J1's total duration is 5 and provides a
lower bound, so the schedule achieves the optimum. In this case,
\texttt{jobSequence} reflects the operation dispatch order and may repeat
names; it must not be interpreted as a permutation of jobs.
Branch and bound with dominance pruning supports up to six jobs and
36 operations. These are educational limits: search time can grow
substantially well before they are reached. Native is deterministic and
marks both variants \texttt{fixtureScale}; Validate does not construct
schedules, and External is unavailable.

\section{Reading the Gantt chart and using macOS}
Import the JSON envelope, whose \texttt{kind} is \texttt{flowShop} or
\texttt{jobShop}. To enter or change operations, use the structured
Definition form or the model JSON. These variants do not use the CPM/PERT
activity editor. They are different problems, and the data for one do not
replace those for the other.
Validate and run with QSB Native. The solution opens Timeline: each row is
a machine, bars are operations, and gaps show idle time.
The horizontal scale lets you examine long schedules. Operation details
identify the job, machine, start, finish, and duration; JSON also contains
\texttt{idleBefore} and the schedules by machine.

Before accepting the result, check that no machine has overlapping bars,
that each job retains its operation order, and that the maximum finish
equals the makespan. A short schedule does not guarantee that due dates
are met; compare them separately if they matter to the actual decision.
The algorithms run in QSBCore, and the Gantt chart presents the same times
exported by the CLI.

\section{Exercises}
\begin{enumerate}
\item Change J2's second operation in the flow shop from 1 to 5.
      Compare both permutations and verify the result.
\item In the job shop, add an initial availability time of 2 for M2.
      Draw the waiting times and check that the solver does not schedule
      operations before that time.
\item Set J1's due date to 4. Explain why changing this field
      does not turn the problem into tardiness minimization.
\end{enumerate}
